\documentclass[12pt]{article}
\usepackage[utf8]{inputenc}
\usepackage[T1]{fontenc}
\usepackage{hyperref}
\title{Wunumara Language Resources for K-12 Curriculum Development}
\author{Compiled by Claude Sonnet 5.5 (claude-sonnet-5-5) for Andr\'as Kornai}
\date{}
\begin{document}
\maketitle

\section{Introduction}
Wunumara (ISO 639-3: wnn) is an extinct Australian language of the Mayi group in Queensland. Its country runs from the headwaters of the Flinders River east to Richmond. Other names include Wanamara, Maykulan and Unamara. Breen (1981) treats it as part of a dialect chain with Mayi-Thakurti, Ngawun, Mayi-Kulan and Mayi-Yapi. Wikipedia's Ngawun article notes some overlap in the ISO entries for the two languages. Very little teaching material appeared in the searches.

\section{Primary Educational Resources}
No primers, textbooks or school programmes were found in the searches.

\section{Linguistic and Reference Materials}
\begin{itemize}
\item AIATSIS Austlang entry G16.1 Wunumara: \url{https://collection.aiatsis.gov.au/austlang/language/g16.1}
\item Wikipedia on the Wanamara people: \url{https://en.wikipedia.org/wiki/Wanamara}
\item Wikipedia on the related Ngawun language, which also notes the ISO code overlap: \url{https://en.wikipedia.org/wiki/Ngawun_language}
\end{itemize}

\section{Religious and Community Resources}
None were found in the searches.

\section{Digital and Audio Resources}
The Wikitongues page for the code lists no documented content that I could confirm: \url{https://wikitongues.org/languages/wnn/}. A University of Sydney archives record also appeared in the results, but I could not confirm that it relates to Wunumara: \url{https://archives-search.sydney.edu.au/nodes/view/19272}.

\section{Curriculum Development Context}
The language has no known speakers. Any revitalisation work should be led by descendant communities, building on the AIATSIS record and archived wordlists. Breen's work on the Mayi languages is the main scholarly lead mentioned in the results, but I did not see a URL for it.
\end{document}